\documentclass[a4paper,12pt,twoside]{memoir}

% Castellano
\usepackage[spanish,es-tabla]{babel}
\selectlanguage{spanish}
\usepackage[utf8]{inputenc}
\usepackage[T1]{fontenc}
\usepackage{lmodern} % Scalable font
\usepackage{microtype}
\usepackage{placeins}

\RequirePackage{booktabs}
\RequirePackage[table]{xcolor}
\RequirePackage{xtab}
\RequirePackage{multirow}

% Links
\PassOptionsToPackage{hyphens}{url}\usepackage[colorlinks]{hyperref}
\hypersetup{
	allcolors = {red}
}

% Ecuaciones
\usepackage{amsmath}

% Rutas de fichero / paquete
\newcommand{\ruta}[1]{{\sffamily #1}}

% Párrafos
\nonzeroparskip

% Huérfanas y viudas
\widowpenalty100000
\clubpenalty100000

% Imagenes
\usepackage{graphicx}
\newcommand{\imagen}[2]{
	\begin{figure}[!h]
		\centering
		\includegraphics[width=0.9\textwidth]{#1}
		\caption{#2}\label{fig:#1}
	\end{figure}
	\FloatBarrier
}

\newcommand{\imagenflotante}[2]{
	\begin{figure}%[!h]
		\centering
		\includegraphics[width=0.9\textwidth]{#1}
		\caption{#2}\label{fig:#1}
	\end{figure}
}



% El comando \figura nos permite insertar figuras comodamente, y utilizando
% siempre el mismo formato. Los parametros son:
% 1 -> Porcentaje del ancho de página que ocupará la figura (de 0 a 1)
% 2 --> Fichero de la imagen
% 3 --> Texto a pie de imagen
% 4 --> Etiqueta (label) para referencias
% 5 --> Opciones que queramos pasarle al \includegraphics
% 6 --> Opciones de posicionamiento a pasarle a \begin{figure}
\newcommand{\figuraConPosicion}[6]{%
  \setlength{\anchoFloat}{#1\textwidth}%
  \addtolength{\anchoFloat}{-4\fboxsep}%
  \setlength{\anchoFigura}{\anchoFloat}%
  \begin{figure}[#6]
    \begin{center}%
      \Ovalbox{%
        \begin{minipage}{\anchoFloat}%
          \begin{center}%
            \includegraphics[width=\anchoFigura,#5]{#2}%
            \caption{#3}%
            \label{#4}%
          \end{center}%
        \end{minipage}
      }%
    \end{center}%
  \end{figure}%
}

%
% Comando para incluir imágenes en formato apaisado (sin marco).
\newcommand{\figuraApaisadaSinMarco}[5]{%
  \begin{figure}%
    \begin{center}%
    \includegraphics[angle=90,height=#1\textheight,#5]{#2}%
    \caption{#3}%
    \label{#4}%
    \end{center}%
  \end{figure}%
}
% Para las tablas
\newcommand{\otoprule}{\midrule [\heavyrulewidth]}
%
% Nuevo comando para tablas pequeñas (menos de una página).
\newcommand{\tablaSmall}[5]{%
 \begin{table}
  \begin{center}
   \rowcolors {2}{gray!35}{}
   \begin{tabular}{#2}
    \toprule
    #4
    \otoprule
    #5
    \bottomrule
   \end{tabular}
   \caption{#1}
   \label{tabla:#3}
  \end{center}
 \end{table}
}

%
% Nuevo comando para tablas pequeñas (menos de una página).
\newcommand{\tablaSmallSinColores}[5]{%
 \begin{table}[H]
  \begin{center}
   \begin{tabular}{#2}
    \toprule
    #4
    \otoprule
    #5
    \bottomrule
   \end{tabular}
   \caption{#1}
   \label{tabla:#3}
  \end{center}
 \end{table}
}

\newcommand{\tablaApaisadaSmall}[5]{%
\begin{landscape}
  \begin{table}
   \begin{center}
    \rowcolors {2}{gray!35}{}
    \begin{tabular}{#2}
     \toprule
     #4
     \otoprule
     #5
     \bottomrule
    \end{tabular}
    \caption{#1}
    \label{tabla:#3}
   \end{center}
  \end{table}
\end{landscape}
}

%
% Nuevo comando para tablas grandes con cabecera y filas alternas coloreadas en gris.
\newcommand{\tabla}[6]{%
  \begin{center}
    \tablefirsthead{
      \toprule
      #5
      \otoprule
    }
    \tablehead{
      \multicolumn{#3}{l}{\small\sl continúa desde la página anterior}\\
      \toprule
      #5
      \otoprule
    }
    \tabletail{
      \hline
      \multicolumn{#3}{r}{\small\sl continúa en la página siguiente}\\
    }
    \tablelasttail{
      \hline
    }
    \bottomcaption{#1}
    \rowcolors {2}{gray!35}{}
    \begin{xtabular}{#2}
      #6
      \bottomrule
    \end{xtabular}
    \label{tabla:#4}
  \end{center}
}

%
% Nuevo comando para tablas grandes con cabecera.
\newcommand{\tablaSinColores}[6]{%
  \begin{center}
    \tablefirsthead{
      \toprule
      #5
      \otoprule
    }
    \tablehead{
      \multicolumn{#3}{l}{\small\sl continúa desde la página anterior}\\
      \toprule
      #5
      \otoprule
    }
    \tabletail{
      \hline
      \multicolumn{#3}{r}{\small\sl continúa en la página siguiente}\\
    }
    \tablelasttail{
      \hline
    }
    \bottomcaption{#1}
    \begin{xtabular}{#2}
      #6
      \bottomrule
    \end{xtabular}
    \label{tabla:#4}
  \end{center}
}

%
% Nuevo comando para tablas grandes sin cabecera.
\newcommand{\tablaSinCabecera}[5]{%
  \begin{center}
    \tablefirsthead{
      \toprule
    }
    \tablehead{
      \multicolumn{#3}{l}{\small\sl continúa desde la página anterior}\\
      \hline
    }
    \tabletail{
      \hline
      \multicolumn{#3}{r}{\small\sl continúa en la página siguiente}\\
    }
    \tablelasttail{
      \hline
    }
    \bottomcaption{#1}
  \begin{xtabular}{#2}
    #5
   \bottomrule
  \end{xtabular}
  \label{tabla:#4}
  \end{center}
}



\definecolor{cgoLight}{HTML}{EEEEEE}
\definecolor{cgoExtralight}{HTML}{FFFFFF}

%
% Nuevo comando para tablas grandes sin cabecera.
\newcommand{\tablaSinCabeceraConBandas}[5]{%
  \begin{center}
    \tablefirsthead{
      \toprule
    }
    \tablehead{
      \multicolumn{#3}{l}{\small\sl continúa desde la página anterior}\\
      \hline
    }
    \tabletail{
      \hline
      \multicolumn{#3}{r}{\small\sl continúa en la página siguiente}\\
    }
    \tablelasttail{
      \hline
    }
    \bottomcaption{#1}
    \rowcolors[]{1}{cgoExtralight}{cgoLight}

  \begin{xtabular}{#2}
    #5
   \bottomrule
  \end{xtabular}
  \label{tabla:#4}
  \end{center}
}


















\graphicspath{ {./img/} }

% Capítulos
\chapterstyle{bianchi}
\newcommand{\capitulo}[2]{
	\setcounter{chapter}{#1}
	\setcounter{section}{0}
	\chapter*{#2}
	\addcontentsline{toc}{chapter}{#2}
	\markboth{#2}{#2}
}

% Apéndices
\renewcommand{\appendixname}{Apéndice}
\renewcommand*\cftappendixname{\appendixname}

\newcommand{\apendice}[1]{
	%\renewcommand{\thechapter}{A}
	\chapter{#1}
}

\renewcommand*\cftappendixname{\appendixname\ }

% Formato de portada
\makeatletter
\usepackage{xcolor}
\newcommand{\tutor}[1]{\def\@tutor{#1}}
\newcommand{\course}[1]{\def\@course{#1}}
\definecolor{cpardoBox}{HTML}{E6E6FF}
\def\maketitle{
  \null
  \thispagestyle{empty}
  % Cabecera ----------------
\noindent\includegraphics[width=\textwidth]{cabecera}\vspace{1cm}%
  \vfill
  % Título proyecto y escudo informática ----------------
  \colorbox{cpardoBox}{%
    \begin{minipage}{.8\textwidth}
      \vspace{.5cm}\Large
      \begin{center}
      \textbf{TFG del Grado en Ingeniería Informática}\vspace{.6cm}\\
      \textbf{\LARGE\@title{}}
      \end{center}
      \vspace{.2cm}
    \end{minipage}

  }%
  \hfill\begin{minipage}{.20\textwidth}
    \includegraphics[width=\textwidth]{escudoInfor}
  \end{minipage}
  \vfill
  % Datos de alumno, curso y tutores ------------------
  \begin{center}%
  {%
    \noindent\LARGE
    Presentado por \@author{}\\ 
    en Universidad de Burgos --- \@date{}\\
    Tutor: \@tutor{}\\
  }%
  \end{center}%
  \null
  \cleardoublepage
  }
\makeatother

\newcommand{\nombre}{Pablo Fernández Bravo} %%% cambio de comando

% Datos de portada
\title{título del TFG}
\author{\nombre}
\tutor{Dr. Carlos López Nozal y \\ Dr. Jesús María Alonso Abad}
\date{\today}

\begin{document}

\maketitle


\newpage\null\thispagestyle{empty}\newpage


%%%%%%%%%%%%%%%%%%%%%%%%%%%%%%%%%%%%%%%%%%%%%%%%%%%%%%%%%%%%%%%%%%%%%%%%%%%%%%%%%%%%%%%%
\thispagestyle{empty}


\noindent\includegraphics[width=\textwidth]{cabecera}\vspace{1cm}

\noindent D. Carlos López Nozal y D. Jesús María Alonso Abad, profesores del departamento de Ingeniería Informática, área de Lenguajes y Sistemas Informáticos.

\noindent Expone:

\noindent Que el alumno D. \nombre, con DNI 71364831V, ha realizado el Trabajo final de Grado en Ingeniería Informática titulado título de TFG. 

\noindent Y que dicho trabajo ha sido realizado por el alumno bajo la dirección del que suscribe, en virtud de lo cual se autoriza su presentación y defensa.

\begin{center} %\large
En Burgos, {\large \today}
\end{center}

\vfill\vfill\vfill

% Author and supervisor
\begin{minipage}{0.45\textwidth}
\begin{flushleft} %\large
Vº. Bº. del Tutor:\\[2cm]
D. Carlos López Nozal
\end{flushleft}
\end{minipage}
\hfill
\begin{minipage}{0.45\textwidth}
\begin{flushleft} %\large
Vº. Bº. del co-tutor:\\[2cm]
D. Jesús María Alonso Abad
\end{flushleft}
\end{minipage}
\hfill

\vfill

% para casos con solo un tutor comentar lo anterior
% y descomentar lo siguiente
%Vº. Bº. del Tutor:\\[2cm]
%D. nombre tutor


\newpage\null\thispagestyle{empty}\newpage




\frontmatter

% Abstract en castellano
\renewcommand*\abstractname{Resumen}
\begin{abstract}
En este primer apartado se hace una \textbf{breve} presentación del tema que se aborda en el proyecto.
\end{abstract}

\renewcommand*\abstractname{Descriptores}
\begin{abstract}
minería de textos, microservicios, github, python, pygithub
\end{abstract}

\clearpage

% Abstract en inglés
\renewcommand*\abstractname{Abstract}
\begin{abstract}
A \textbf{brief} presentation of the topic addressed in the project.
\end{abstract}

\renewcommand*\abstractname{Keywords}
\begin{abstract}
text mining, microservices, github, python, pygithub
\end{abstract}

\clearpage

% Indices
\tableofcontents

\clearpage

\listoffigures

\clearpage

\listoftables
\clearpage

\mainmatter
\capitulo{1}{Introducción}

Los proyectos software requieren del establecimiento de una serie de medidas que permitan organizar y gestionar su progreso a lo largo de las diversas fases del desarrollo. Entre estas medidas en los proyectos software destacan los \textbf{sistemas de gestión y seguimiento de incidencias}. La tendencia en el uso de metodologías rápidas en este tipo de proyectos implica la necesidad de dividir los grandes proyectos en cortas iteraciones compuestas por pequeñas tareas. La asignación de este tipo de tareas suele llevarse a cabo de manera individual, debido a que de manera diaria se realizan reuniones donde exponer los problemas encontrados y proponer soluciones e ideas en conjunto.

La organización de este tipo de reuniones no siempre tiene porqué llevarse a cabo de manera presencial, es posible realizarse a través de foros internos donde debatir las propuestas, clasificar las tareas, gestionar las incidencias detectadas, y en definitiva analizar la situación de las pequeñas tareas con el objetivo de concluir el proyecto de acuerdo con los requisitos en tiempo y recursos establecidos. En este punto es donde surgen programas o plataformas con la finalidad de simplificar la gestión de las tareas en grandes proyectos, cuya organización puede resultar compleja y caótica.

El inconveniente de este tipo de plataformas recae en que pese a facilitar las labores de seguimiento del estado de los procesos o incidencias, siguen requiriendo un alto grado de atención manual por parte del encargado de su gestión. La razón detrás de ello recae en el hecho de que la definición, descripción, debate, y conclusión de las tareas poseen un alto componente social que no podía ser interpretado por las máquinas.

El punto de inflexión con los problemas anteriores llega en los últimos 10 años, momento en el cual las empresas comenzaron a descubrir el potencial de la inmensa cantidad de información que estas habían almacenado durante décadas. La necesidad de llevar a cabo un tratamiento de estos datos propició el creciente desarrollo y evolución de las \textbf{técnicas de procesamiento del lenguaje natural (PLN)}. Estas técnicas tienen como principal objetivo la aplicación de \textbf{Inteligencia Artificial (IA)} para el análisis automático de las comunicaciones humanas por medio de los computadores, permitiendo la extracción, clasificación y descubrimiento de detalles e información relevante que no aparece dispuesta de manera explícita.

Estas técnicas de procesamiento del lenguaje natural en conjunto con la necesidad de automatizar aquellas labores de seguimiento de incidencias en la gestión de un proyecto es el lugar donde se sitúa nuestro objetivo con este trabajo. Con la finalidad de comprobar y demostrar el potencial de esta especialidad del aprendizaje automático en su aplicación a las labores de seguimiento de incidencias de un proyecto software este trabajo plantea el diseño de una plataforma con la que experimentar con distintos modelos de PLN aplicables y su ajuste a través de la modificación de parámetros.  

Los modelos de PLN que se ha decidido incluir en la plataforma permiten al usuario conocer la probabilidad con la que el contenido de una incidencia se ajusta a una serie de etiquetas definidas tanto en el propio repositorio como introducidas por el usuario (\textbf{\emph{Zero-Shot Classification}}), generar resúmenes abstractivos que sinteticen el contenido de las incidencias (\textbf{\emph{Summarization}}), y obtener gráficos que señale la postura de los participantes con respecto a una incidencia deducida a partir del análisis de su actitud en los comentarios (\textbf{\emph{Sentiment Analysis}}).

\capitulo{2}{Objetivos del proyecto}

Este capítulo tiene como objeto el recoger los objetivos generales que se establecieron al comienzo del proyecto con la finalidad de determinar y acotar el alcance del proyecto, así como los objetivos técnicos con el propósito de determinar las técnicas y tecnologías clave que se acordó utilizar.

\section{Objetivos generales}

En este apartado se exponen los objetivos de carácter general que se pactó implementar en el proyecto. A continuación se describen estos objetivos:

\begin{itemize} [\textbullet]

\item Realizar una aplicación web que permita al usuario experimentar con múltiples técnicas de procesamiento del lenguaje natural (PLN). Tras estudiar los posibles métodos que aplicar se decidió que en esta ocasión la plataforma trabajará con tareas de \emph{clasificación Zero-Shot}, el análisis de sentimientos y la generación de resúmenes abstractivos.

\item Los modelos que se van a utilizar no se implementarán y entrenarán desde cero, en su lugar se recurrirá a la explotación del potencial de los \textbf{modelos preentrenados}. Mediante portales como \emph{Hugging Face} \cite{platform:hugging_face} es posible obtener modelos que aplican en multitud de campos con el beneficio de haber sido entrenados y ajustados por anterioridad. Este tipo de iniciativas realza el valor del software de código abierto y permiten que pequeños proyectos dispongan de acceso a potentes herramientas que ofrecen resultados con suficiente grado de calidad.

\item Los datos sobre los que se aplicarán los modelos deberán ser extraídos desde las incidencias de repositorios de código abierto de \emph{GitHub} \cite{platform:github}. La motivación de que los datos de entrada se obtengan mediante esta vía surge con el fin de demostrar los beneficios de aplicar este tipo de metodologías a la gestión de tareas dentro de un proyecto. A partir de la aplicación se podría realizar una clasificación automática o semiautomática de tareas, obtener un análisis del estado moral del equipo u ofrecer una sinopsis de las tareas y su estado.

\item Enfocar la \hyperref[sec:arquitectura]{arquitectura y diseño del proyecto} teniendo en cuenta la posibilidad de aplicar mejoras sobre el mismo y ampliar sus funcionalidades de cara futuras iteraciones tanto en el \emph{front-end} como en el \emph{back-end}. Esto incluye dotar al código de una buena documentación que permita a futuros desarrolladores comprender el funcionamiento interno de la aplicación.
\item Trasladar a la práctica aquellos conocimientos adquiridos a lo largo del grado en las materias de Ingeniería del Software, Gestión de Proyectos, Metodología de la Programación, Bases de Datos, Algoritmia, Diseño y Mantenimiento del Software, y Minería de Datos.
\item Ofrecer el código bajo una licencia que posibilite su libre uso y acceso, permitiendo la continuidad del proyecto más allá del alcance de este trabajo.

\end{itemize}

\bigskip

\section{Objetivos técnicos}

Este apartado tiene como finalidad presentar aquellos desafíos técnicos que afrontar de cara a dotar al proyecto de la tecnología necesaria para la implementación de sus funciones. Estos objetivos mencionados son los siguientes:

\begin{itemize} [\textbullet]

\item El proyecto debe plantearse teniendo en mente una \hyperref[sec:microservicios]{\textbf{arquitectura de microservicios}} que otorgue de independencia a cada fase del proceso. El objetivo de estar arquitectura consiste en otorgar la libertad de que en caso de que se requiriese una ampliación o reubicación de recursos por parte de uno de los servicios, estos puedan ser ajustados de manera individual sin alterar el funcionamiento completo de la aplicación. 

\item Se deberá contar con un servicio encargado de establecer la conexión con la \emph{API} de \emph{GitHub} \cite{platform:github_api} para la \hyperref[sec:extraccion]{extracción de los datos} Este servicio a su vez se encargará de extraer aquella información que así se requiera de las incidencias, aplicar un preprocesado sobre estos datos, y una vez se obtengan los datos en limpio, se proceda a su almacenamiento en la base de datos.

\item La base de datos que se escoja deberá soportar los accesos de manera concurrente por parte de los múltiples servicios que compondrán la aplicación. Esta se desplegará en un servicio independiente para facilitar su mantenimiento.

\item Los \hyperref[sec:preentrenados]{\textbf{modelos preentrenados}} ofrecen la posibilidad de modificar ciertos parámetros con el objetivo de adaptar los resultados a la salida deseada. Siempre que sea posible se ofrecerá al usuario de la aplicación la posibilidad de configurar estos parámetros de acuerdo con sus preferencias para dotar a la aplicación de una mayor funcionalidad.

\item Con el fin de mejorar los resultados de los \hyperref[sec:preentrenados]{\textbf{modelos preentrenados}} se estudiarán las limitaciones en el tamaño de sus entradas y se implementarán técnicas de \hyperref[sec:preprocesamiento]{preprocesado} que permitan mitigar sus efectos en situaciones donde los textos sean demasiado grandes.

\end{itemize}

\capitulo{3}{Conceptos teóricos}

En aquellos proyectos que necesiten para su comprensión y desarrollo de unos conceptos teóricos de una determinada materia o de un determinado dominio de conocimiento, debe existir un apartado que sintetice dichos conceptos.

Algunos conceptos teóricos de \LaTeX \footnote{Créditos a los proyectos de Álvaro López Cantero: Configurador de Presupuestos y Roberto Izquierdo Amo: PLQuiz}.

\section{Secciones}

Las secciones se incluyen con el comando section.

\subsection{Subsecciones}

Además de secciones tenemos subsecciones.

\subsubsection{Subsubsecciones}

Y subsecciones. 


\section{Referencias}

Las referencias se incluyen en el texto usando cite \cite{wiki:latex}. Para citar webs, artículos o libros \cite{koza92}.


\section{Imágenes}

Se pueden incluir imágenes con los comandos standard de \LaTeX, pero esta plantilla dispone de comandos propios como por ejemplo el siguiente:

\imagen{escudoInfor}{Autómata para una expresión vacía}



\section{Listas de items}

Existen tres posibilidades:

\begin{itemize}
	\item primer item.
	\item segundo item.
\end{itemize}

\begin{enumerate}
	\item primer item.
	\item segundo item.
\end{enumerate}

\begin{description}
	\item[Primer item] más información sobre el primer item.
	\item[Segundo item] más información sobre el segundo item.
\end{description}
	
\begin{itemize}
\item 
\end{itemize}

\section{Tablas}

Igualmente se pueden usar los comandos específicos de \LaTeX o bien usar alguno de los comandos de la plantilla.

\tablaSmall{Herramientas y tecnologías utilizadas en cada parte del proyecto}{l c c c c}{herramientasportipodeuso}
{ \multicolumn{1}{l}{Herramientas} & App AngularJS & API REST & BD & Memoria \\}{ 
HTML5 & X & & &\\
CSS3 & X & & &\\
BOOTSTRAP & X & & &\\
JavaScript & X & & &\\
AngularJS & X & & &\\
Bower & X & & &\\
PHP & & X & &\\
Karma + Jasmine & X & & &\\
Slim framework & & X & &\\
Idiorm & & X & &\\
Composer & & X & &\\
JSON & X & X & &\\
PhpStorm & X & X & &\\
MySQL & & & X &\\
PhpMyAdmin & & & X &\\
Git + BitBucket & X & X & X & X\\
Mik\TeX{} & & & & X\\
\TeX{}Maker & & & & X\\
Astah & & & & X\\
Balsamiq Mockups & X & & &\\
VersionOne & X & X & X & X\\
} 

\capitulo{4}{Técnicas y herramientas}

Esta parte de la memoria tiene como objetivo presentar las técnicas metodológicas y las herramientas de desarrollo que se han utilizado para llevar a cabo el proyecto. Si se han estudiado diferentes alternativas de metodologías, herramientas, bibliotecas se puede hacer un resumen de los aspectos más destacados de cada alternativa, incluyendo comparativas entre las distintas opciones y una justificación de las elecciones realizadas. 
No se pretende que este apartado se convierta en un capítulo de un libro dedicado a cada una de las alternativas, sino comentar los aspectos más destacados de cada opción, con un repaso somero a los fundamentos esenciales y referencias bibliográficas para que el lector pueda ampliar su conocimiento sobre el tema.



\capitulo{5}{Aspectos relevantes del desarrollo del proyecto}

\section{Arquitectura del proyecto} \label{sec:arquitectura}

\subsection{Arquitectura de Microservicios} \label{sec:microservicios}
La planificación de la arquitectura sobre la cual se iba a desarrollar el proyecto quedó definida en las primeras fases de exploración de la idea. Su elección estuvo condicionada al cumplimiento de una serie de características que se planteaban como clave en el desarrollo de una plataforma moderna. Entre las propiedades de las que se pretendía dotar a la aplicación destacan las siguientes:

\begin{itemize}
    \item [\textbullet] \textbf{Bajo acoplamiento}. Se decidió buscar la separación entre sí de aquellos componentes que no dependiesen del resto para su ejecución. Una de las principales ventajas de este tipo de planteamientos reside en minimizar el impacto que puede tener una falla que se de en alguna de las fases del procesamiento que se llevan a cabo en el back-end. Otra de sus ventajas implica en proporcionar la posibilidad de optimizar los recursos a las necesidades de cada fase, asignando una menor cantidad de recursos para aquellas tareas de bajo coste computacional, y distribuyendo esos recursos en aquellos procesos que requieran de una mayor potencia de cálculo y/o almacenamiento.
    \item [\textbullet] \textbf{Mantenible y adaptable a cambios}. La esencia del proyecto resalta el concepto de experimentar con la aplicación de modelos de PLN a los sistemas de seguimiento de incidencia con el objetivo de lograr simplificar y mejorar partes del proceso. Esta propuesta deja abierta la posibilidad de incorporar nuevas tareas y funcionalidades a la plataforma de cara a un futuro, así como la actualización de modelos a nuevas implementaciones.
    \item [\textbullet] \textbf{Alta disponibilidad}. La extracción de los datos y la aplicación de los modelos de PLN son procesos que demandan de una cantidad moderada de tiempo para la realización de sus funciones. Estos retrasos no deben repercutir en la experiencia del usuario, al cual se le debe seguir permitiendo interactuar con la plataforma mientras estos procesos se ejecutan en segundo plano.
\end{itemize}

Con el afán de cumplir con la implementación de una arquitectura que cumpliese dotase a la plataforma de las características anteriormente mencionadas nos decantamos por una arquitectura basada en microservicios. Este patrón arquitectónico permite encapsular las diferentes tareas de la aplicación en servicios independientes con un alto nivel de especificidad.

Una de las ventajas principales reside en el hecho de que la separación del servicio API REST encargado de gestionar las peticiones desde la aplicación, de los servicios de extracción o procesado, permite atender y encolar peticiones mientras se atienden peticiones previas evitando generar tiempos muertos. La segunda virtud de este tipo de arquitectura otorga a los mantenedores de la plataforma la opción de duplicar servicios en caso de que la demanda de este tipo de tareas supere a la capacidad de procesamiento de los servicios actuales.

\subsection{Arquitectura dirigida por eventos} \label{sec:dirigida_por_eventos}

La decisión arquitectónica escogida entorno a la separación de los diferentes servicios en componentes trae consigo la necesidad de implementar algún mecanismo mediante el cual los componentes sean capaces de detectar cuando deben llevar a cabo sus labores. El mecanismo requerido nos condujo directamente a la implantación de una arquitectura basada en eventos.
Este patrón arquitectónico suele aplicarse en entornos asíncronos, es decir, que cuando se genera una petición esta no siempre va a poder ser atendida y resuelta en ese mismo instante. Las peticiones entrantes son atendidas por un agente, este se encarga de emitir un evento y redireccionarlo a la cola de tareas del correspondiente servicio. Por su parte, los servicios tienen la responsabilidad de atender a dichos eventos y actuar según proceda manteniendo un orden de acuerdo con las prioridades establecidas. Seguidamente se presentan los detalles de la implantación de este patrón en nuestro proyecto:

\vspace{-0.5cm}
\begin{itemize}
\item [\textbullet] Agente generador de eventos. El agente generador de eventos de nuestra plataforma es el servicio de API REST, el cual se encarga de atender las peticiones lanzadas por los usuarios a través de la web. Las peticiones son procesadas, retornando una respuesta informando de si esta ha sido aceptada o rechazada debido a algún tipo de error en los parámetros requeridos.  Ante una petición correcta el sistema genera una nueva entrada en la base de datos detallando el tipo de tarea que ha sido solicitada, y los atributos asociados a dicha tarea.
\item [\textbullet] Canal de comunicación de eventos. La implantación del canal a través del cual los servicios son conscientes de las peticiones entrantes es una tabla de la base de datos donde cada fila representa una nueva tarea asociada a un identificador único. En los atributos de cada entrada es donde se indica hacia qué servicio va dirigida la petición, una bandera que indica el estado en el que se encuentra (esperando, en proceso, resuelta, fallida o desactualizada).
\item [\textbullet] Consumidores de eventos. Los servicios consumidores de eventos son tanto el servicio de procesamiento como el de extracción. Estos servicios acceden de manera periódica a la base de datos en busca de nuevas peticiones que requieran de su atención. En el momento se detecta un nuevo evento estos recuperan los argumentos de la petición y actualizan su estado. Finalmente, una vez resuelta la petición actualizan el estado de la tarea de acorde a su resolución.
\end{itemize}

El alto grado sobre el cual se ha implementado la aplicación no se encuentra libre de inconvenientes. Entre ellos debemos destacar la necesidad del agente externo que lanza las solicitudes contra la API REST de encargarse también de recoger los datos consultando a dicho servicio actualizaciones sobre el estado en el que se encuentra su petición. Este efecto se ve mitigado debido a que será el servidor que maneja la aplicación web el encargado de realizar llamadas de manera periódica hasta que vea resuelta su petición.


\capitulo{6}{Trabajos relacionados}

Este apartado sería parecido a un estado del arte de una tesis o tesina. En un trabajo final grado no parece obligada su presencia, aunque se puede dejar a juicio del tutor el incluir un pequeño resumen comentado de los trabajos y proyectos ya realizados en el campo del proyecto en curso. 

\capitulo{7}{Conclusiones y Líneas de trabajo futuras}

\section{Conclusiones}

El proyecto \textbf{GitHub Text Mining} es el resultado del esfuerzo invertido durante un plazo de aproximadamente \textbf{6 meses} con el ambicioso objetivo de lograr desarrollar todos los aspectos que hacen a una plataforma funcional desde cero. En este trabajo se llevado a cabo un desarrollo \textbf{full stack}, o lo que es lo mismo, se ha desarrollado tanto el aspecto interno de una aplicación, trabajando con \textbf{tecnologías de servidor}, como el desarrollo de aquellos aspectos relacionados con la \textbf{presentación de la información} y el \textbf{establecimiento de conexiones con la fuente de los datos}.

Las principal conclusión que se extrae del desarrollo de este proyecto es \textbf{la importancia de tomar una decisión y ser capaz de defenderla}. En un trabajo de tal magnitud es inevitable enfrentarse a momentos en los que afrontar decisiones que marcarán la dirección del proyecto a futuro. Estas decisiones son tomadas en base a los conocimientos que se poseen en dicho momento, y por lo tanto, es común que con el paso del tiempo y la adquisición de experiencia se formulen nuevas propuestas con la finalidad de perfeccionar las implementaciones en apartados previos. Es necesario establecerse límites en cuanto a \textbf{asumir que un proyecto nunca va a estar completo}, y consecuentemente ceñirse en sacar a delante los objetivos propuestos antes de tratar de actualizar aspectos que actualmente cumplen con su función.

La segunda conclusión extraída de este proyecto, directamente relacionada con la mencionada en el anterior párrafo, es \textbf{la importancia de invertir tiempo en estudiar la dirección que va a tomar un proyecto y las alternativas disponibles} antes de aferrarse por completo a una idea. En un proyecto de gran magnitud, invertir recursos en la investigación y desarrollo de una buena arquitectura que dote a un proyecto de robustez resulta un aspecto clave de cara a la incorporación de cambios y novedades en un futuro.

Para finalizar, se destaca la relevancia de afrontar este tipo de Trabajos de Fin de Grado como una oportunidad para \textbf{aplicar los conocimientos adquiridos} a lo largo del Grado, y aprovechar para indagar sobre aquellas áreas de interés en las que tal vez no se haya profundizado tanto como a uno le gustaría. \textbf{Aprovechar la libertad creativa del proyecto} es una forma de descubrir que facetas de la informática encajan con nuestra perspectiva y establecer nuestro próximo objetivo personal.

\section{Líneas de trabajo futuras}

En esta sección se incluyen ideas acerca de la incorporación de nuevas funcionalidades que quedaron descartadas debido al alcance establecido para el proyecto, así como posibles áreas de mejora sobre las que seguir trabajando.

\subsection{Mejoras e incorporación de nuevas tareas}

Desde el inicio el proyecto GitHub Text Mining se planteó con la idea de establecer una arquitectura base que permitiese la ampliación de los modelos soportados a otro tipo de tareas. Actualmente los experimentos que se incluyen permiten obtener una clasificación de las incidencias entorno a etiquetas, generar resúmenes a partir su contenido y realizar un análisis de sentimientos sobre las incidencias, incluyendo filtros de usuarios y comentarios. En adición a la introducción de mejoras en las respectivas implementaciones de estas tareas, se plantea la introducción de nuevas funcionalidades que permitan detectar incidencias duplicadas en un repositorio haciendo uso de modelos de \textbf{Sentence Similarity}, así como estudiar la introducción de tareas de \textbf{Traducción}.

\subsection{Implementar acceso mediante el uso de tokenes de acceso personal por usuario}

Una decisión de diseño tomada inicialmente condiciona a la utilización de un único token de acceso personal para el establecimiento de la conexión con la API de GitHub. Este factor puede resultar limitante en caso de que se pretenda desplegar la aplicación a causa de las \textbf{limitaciones de la API de GitHub} \cite{lf:github_api_limits}. Sería preferible diseñar alguna estrategia alternativa en la que cada usuario deba proporcionar su propio token para poder utilizar la plataforma. Se deberá tener en cuenta que estos tokenes son información personal, y por lo tanto deberán transferirse, tratarse y almacenarse de forma segura.

\subsection{Refractorización de la API REST}

La \emph{API REST} ha ido creciendo conforme se requerían nuevas formas de acceder a los datos desde la aplicación web. Esto ha propiciado la inclusión de una gran cantidad de métodos que deberían ser refractorizados con el objetivo de reducir su complejidad y deshacerse de aquellos \emph{endpoints} que resulten prescindibles. Otro punto de mejora a considerar implicaría la inclusión de un sistema que estableciese limitaciones en el número de peticiones que se permiten recibir desde una misma dirección durante un cierto periodo de tiempo.



\bibliographystyle{plain}
\bibliography{bibliografia}

\end{document}
